\documentclass[a4paper,12pt]{article}
\usepackage[utf8]{inputenc}
\usepackage[ngerman]{babel}
\usepackage{amsmath,amssymb}
\usepackage{geometry}
\geometry{margin=2.5cm}

\title{Aufgabe 7(c) -- Gauß- und Gauß-Jordan-Elimination}
\author{}
\date{}

\begin{document}
\maketitle

\section*{Gleichungssystem}

\[
\begin{cases}
x + 2y + t + w = 5 \\
-z + 2t - 3w = 3 \\
2x + 4y + z + 5w = 7
\end{cases}
\]

Variablen: $x, y, z, t, w$ (5 Unbekannte, 3 Gleichungen).

\section*{Erweiterte Koeffizientenmatrix}

\[
\left(\begin{array}{ccccc|c}
1 & 2 & 0 & 1 & 1 & 5 \\
0 & 0 & -1 & 2 & -3 & 3 \\
2 & 4 & 1 & 0 & 5 & 7
\end{array}\right)
\]

\section*{Gauß-Elimination (Stufenform)}

\textbf{Schritt 1:} $Z_3 \leftarrow Z_3 - 2Z_1$

\[
\left(\begin{array}{ccccc|c}
1 & 2 & 0 & 1 & 1 & 5 \\
0 & 0 & -1 & 2 & -3 & 3 \\
0 & 0 & 1 & -2 & 3 & -3
\end{array}\right)
\]

\textbf{Schritt 2:} $Z_3 \leftarrow Z_3 + Z_2$

\[
\left(\begin{array}{ccccc|c}
1 & 2 & 0 & 1 & 1 & 5 \\
0 & 0 & -1 & 2 & -3 & 3 \\
0 & 0 & 0 & 0 & 0 & 0
\end{array}\right)
\]

\textbf{Ergebnis:} 2 Pivot-Spalten (Spalte 1 und Spalte 3), keine Widerspruchszeile.

Das System ist \textbf{konsistent} mit \textbf{unendlich vielen Lösungen}.

Freie Variablen: $y, t, w$ (3 freie Variable).

\subsection*{Rückwärtseinsetzen}

Aus Zeile 2: $-z + 2t - 3w = 3 \Rightarrow z = 2t - 3w - 3$

Aus Zeile 1: $x + 2y + t + w = 5 \Rightarrow x = 5 - 2y - t - w$

\section*{Gauß-Jordan-Elimination (reduzierte Stufenform)}

\textbf{Bemerkung:} Die Gauß-Jordan-Elimination bringt hier kaum etwas Neues.
Das Ziel von Gauß-Jordan ist es, auch \textbf{oberhalb} der Pivots zu eliminieren, damit jede
Pivot-Variable nur noch in \textbf{einer} Gleichung vorkommt.

Aber in unserem Fall steht in Zeile 1 bereits eine $0$ in der Pivot-Spalte von Zeile 2 (Spalte 3).
Es gibt also \textbf{nichts zu eliminieren} -- wir müssen nur noch das Pivot in Zeile 2 normieren.

\textbf{Schritt 3:} $Z_2 \leftarrow -Z_2$ (Pivot normieren)

\[
\left(\begin{array}{ccccc|c}
1 & 2 & 0 & 1 & 1 & 5 \\
0 & 0 & 1 & -2 & 3 & -3 \\
0 & 0 & 0 & 0 & 0 & 0
\end{array}\right)
\]

Dies ist die \textbf{reduzierte Zeilenstufenform} (RREF). Man liest ab:
\begin{align*}
\text{Zeile 1:} \quad & x + 2y + t + w = 5 \\
\text{Zeile 2:} \quad & z - 2t + 3w = -3
\end{align*}

Das ist \textbf{exakt dasselbe} wie beim Rückwärtseinsetzen oben.
Die Gauß-Jordan-Elimination hat hier also nur das Vorzeichen gedreht -- die Information ist dieselbe.

\newpage
\section*{Allgemeine Lösung -- Schritt für Schritt erklärt}

Ab hier kommt der Teil, der über die reine Elimination hinausgeht.
Wir haben die RREF:

\[
\left(\begin{array}{ccccc|c}
1 & 2 & 0 & 1 & 1 & 5 \\
0 & 0 & 1 & -2 & 3 & -3 \\
0 & 0 & 0 & 0 & 0 & 0
\end{array}\right)
\]

\subsection*{Schritt 1: Pivot-Variablen vs.\ freie Variablen}

\textbf{Pivot-Spalten} sind die Spalten mit einer führenden 1: Spalte 1 ($x$) und Spalte 3 ($z$).

$\Rightarrow$ $x$ und $z$ sind \textbf{Pivot-Variablen}: Sie werden durch die Gleichungen bestimmt
(sobald man die restlichen Variablen kennt).

\textbf{Spalten ohne Pivot}: Spalte 2 ($y$), Spalte 4 ($t$), Spalte 5 ($w$).

$\Rightarrow$ $y$, $t$ und $w$ sind \textbf{freie Variablen}: Für sie gibt es keine eigene Gleichung.
Sie dürfen \textbf{jeden beliebigen Wert} annehmen -- das System ist trotzdem erfüllt.

\subsection*{Schritt 2: Freie Variablen als Parameter}

Wir geben den freien Variablen neue Namen (um sie von den berechneten zu unterscheiden):

\[
y = s, \quad t = r, \quad w = q \qquad \text{mit } s, r, q \in \mathbb{R}
\]

Das bedeutet nur: \glqq $y$ darf irgendeine reelle Zahl sein -- wir nennen sie $s$.\grqq{}
Die Buchstaben $s, r, q$ sind Platzhalter für \textbf{alle möglichen Wahlen}.

\subsection*{Schritt 3: Pivot-Variablen ausdrücken}

Jetzt setzen wir $y = s$, $t = r$, $w = q$ in die zwei Gleichungen ein:

\begin{align*}
\text{Aus Zeile 2:} \quad z &= -3 + 2r - 3q \\
\text{Aus Zeile 1:} \quad x &= 5 - 2s - r - q
\end{align*}

\subsection*{Schritt 4: Closed Form -- alles zusammen hinschreiben}

Alle 5 Variablen in einer Zeile als Tupel:

\[
\boxed{(x,y,z,t,w) = (5 - 2s - r - q,\; s,\; -3 + 2r - 3q,\; r,\; q), \quad s, r, q \in \mathbb{R}}
\]

Jede konkrete Wahl von $s, r, q$ liefert eine konkrete Lösung. Z.\,B.:
\begin{itemize}
\item $s = 0, r = 0, q = 0 \Rightarrow (5, 0, -3, 0, 0)$
\item $s = 1, r = 0, q = 0 \Rightarrow (3, 1, -3, 0, 0)$
\item $s = 0, r = 1, q = 1 \Rightarrow (3, 0, -4, 1, 1)$
\end{itemize}

\subsection*{Schritt 5: Expanded Form -- in Vektoren zerlegen}

Die Idee: Trenne den \textbf{fixen Anteil} (wenn alle Parameter $= 0$ sind)
von den \textbf{veränderbaren Anteilen} (was sich ändert, wenn $s$, $r$ oder $q$ variiert).

Schreibe die Closed Form untereinander und sortiere nach Parametern:

\[
\begin{pmatrix} x \\ y \\ z \\ t \\ w \end{pmatrix}
=
\begin{pmatrix} 5 \\ 0 \\ -3 \\ 0 \\ 0 \end{pmatrix}
+
\begin{pmatrix} -2s \\ s \\ 0 \\ 0 \\ 0 \end{pmatrix}
+
\begin{pmatrix} -r \\ 0 \\ 2r \\ r \\ 0 \end{pmatrix}
+
\begin{pmatrix} -q \\ 0 \\ -3q \\ 0 \\ q \end{pmatrix}
\]

Den Parameter kann man vor den Vektor ziehen:

\[
\begin{pmatrix} x \\ y \\ z \\ t \\ w \end{pmatrix}
=
\underbrace{\begin{pmatrix} 5 \\ 0 \\ -3 \\ 0 \\ 0 \end{pmatrix}}_{v_0}
+ s \underbrace{\begin{pmatrix} -2 \\ 1 \\ 0 \\ 0 \\ 0 \end{pmatrix}}_{v_1}
+ r \underbrace{\begin{pmatrix} -1 \\ 0 \\ 2 \\ 1 \\ 0 \end{pmatrix}}_{v_2}
+ q \underbrace{\begin{pmatrix} -1 \\ 0 \\ -3 \\ 0 \\ 1 \end{pmatrix}}_{v_3}
\]

\[
\boxed{\vec{x} = v_0 + s \cdot v_1 + r \cdot v_2 + q \cdot v_3, \quad s, r, q \in \mathbb{R}}
\]

\subsection*{Was bedeuten die Vektoren?}

\begin{itemize}
\item $v_0 = (5, 0, -3, 0, 0)^T$ ist \textbf{eine} konkrete Lösung (die \glqq Basislösung\grqq{} für $s = r = q = 0$).
\item $v_1 = (-2, 1, 0, 0, 0)^T$ beschreibt, wie sich die Lösung ändert, wenn man $y$ (also $s$) um 1 erhöht.
\item $v_2 = (-1, 0, 2, 1, 0)^T$ beschreibt die Änderung, wenn man $t$ (also $r$) um 1 erhöht.
\item $v_3 = (-1, 0, -3, 0, 1)^T$ beschreibt die Änderung, wenn man $w$ (also $q$) um 1 erhöht.
\end{itemize}

\textbf{Anschaulich:} $v_0$ ist ein Startpunkt. $v_1, v_2, v_3$ sind drei Richtungen,
in die man sich bewegen kann, ohne die Gleichungen zu verletzen.
Jede Kombination aus diesen Richtungen ergibt eine gültige Lösung.

\subsection*{Woher kommen die Vektoren konkret?}

Man liest sie direkt aus der Closed Form ab. Für jede Variable notiert man:
\begin{center}
\begin{tabular}{c|c|cccc}
Variable & fixer Anteil & Koeff.\ von $s$ & Koeff.\ von $r$ & Koeff.\ von $q$ \\ \hline
$x = 5 - 2s - r - q$ & $5$ & $-2$ & $-1$ & $-1$ \\
$y = s$ & $0$ & $1$ & $0$ & $0$ \\
$z = -3 + 2r - 3q$ & $-3$ & $0$ & $2$ & $-3$ \\
$t = r$ & $0$ & $0$ & $1$ & $0$ \\
$w = q$ & $0$ & $0$ & $0$ & $1$ \\
\end{tabular}
\end{center}

Jede \textbf{Spalte} dieser Tabelle ist ein Vektor: Die erste Spalte ist $v_0$, die zweite $v_1$, usw.

\section*{Probe}

Setze z.\,B.\ $s = 0, r = 0, q = 0$, also $(x,y,z,t,w) = (5,0,-3,0,0)$:

\begin{align*}
5 + 2(0) + 0 + 0 &= 5 \checkmark \\
-(-3) + 2(0) - 3(0) &= 3 \checkmark \\
2(5) + 4(0) + (-3) + 5(0) &= 10 - 3 = 7 \checkmark
\end{align*}

\end{document}
